\documentclass[11pt, a4paper]{article}

\usepackage[a4paper, top=2.5cm, bottom=2.5cm, left=2cm, right=2cm]{geometry}
\usepackage{amsmath}
\usepackage{amssymb}
\usepackage[colorlinks=true, linkcolor=blue, urlcolor=blue, citecolor=blue]{hyperref}

\title{Theory of Everything - Volume 28 \\ \Large Evolutionary Algorithms}
\author{Omega Protocol}
\date{\today}

\begin{document}

\maketitle

\section*{Layman's Summary}
Evolution is usually taught as a biological process on Earth. Volume 28 proves that evolution is actually a universal mathematical algorithm. Charles Darwin's theory of natural selection is just a specific biological example of a much deeper thermodynamic law that forces information networks to optimize themselves over time.

\section{Universal Darwinism (Axiom 31)}
We establish that evolution applies to anything that can copy itself with slight variations—whether it's DNA, computer code, or galaxies. Any informational sub-network capable of replication will inevitably undergo natural selection. Why? Because the universe mathematically favors systems that process and dissipate energy most efficiently.

\section{The Fitness Landscape (Theorem)}
Imagine a foggy, mountainous terrain where the highest peaks represent the best possible survival strategies. We prove that evolution is exactly a mathematical "gradient ascent" on this topological landscape. Species blindfoldedly climb the hills of this landscape, pushed upward by the thermodynamic pressure to survive and reproduce.

\section{Punctuated Equilibrium (Theorem)}
The fossil record shows that species stay exactly the same for millions of years, and then suddenly change rapidly. We prove this is mathematically required. Species get stuck in deep "valleys" (local minima) of the fitness landscape where change is punished. Rapid evolution only happens when chaotic environmental shifts force them to "tunnel" to a new valley.

\end{document}
